\properlane{reconciled-creation}{The Creator Restores Communion by Giving Himself}
The King acclaimed in the Mass created the beings over whom he rules. Colossians identifies him as their sustainer and as the one through whose blood they are reconciled. Chrysostom explains why the Creator himself accomplishes this reconciliation through the Cross. Bellarmine distinguishes Christ's authority over all from the willing obedience of those who acknowledge him. This interpretation asks how creatures estranged from God come to share the peace of his kingdom.

\subsection{The firstborn and the origin of every power}
Colossians names the Son the firstborn of all creation. The phrase can appear to place him at the head of a sequence of creatures. Chrysostom refuses that inference by following the clauses that explain it: all things were created through him, and all continue to depend on him. The Son cannot be assigned to the class whose existence he causes. The vocabulary of birth must be read within the passage's account of creative agency.\footnote{Chrysostom, \emph{Homilies on Colossians} III, on Col 1:15--17, NPNF I.13.}

Chrysostom's exposition develops the point through the distinction between visible and invisible creatures. Invisible powers might seem to offer a higher route to God than the human Christ. Paul names them within creation precisely to exclude that exemption. Thrones, dominations, principalities and powers are not rival sources of being. Their greatness remains received. The Son's primacy is therefore more radical than superiority in rank within a common order.

Nor does his work stop with an initial act of creation. Chrysostom explains that creatures would cease to stand if separated from the Son's sustaining providence. Their present existence is dependent. The ruler who resists him, the angel who serves him and the ordinary believer who prays to him share this creaturely dependence, though their moral relations to him differ. This distinction will matter when Bellarmine distinguishes universal authority from willing obedience.

The Epistle next names the Son head of the Church. Chrysostom emphasizes his condescension in becoming the head of those whose existence depends on him. The one called firstborn of creation is also called firstborn from the dead. The Church's members thus receive their new life from the risen Son through whom they were created.

The Preface names the universal King as the priest who offers himself on the Cross. This sacrifice concerns the enmity of creatures who depend on him for their existence. The distinction between creation and reconciliation matters here: they already belong to him as creatures, and through his death they can come to live with him in peace.

\subsection{Why the blood of the Cross matters to the argument}
Chrysostom distinguishes the greatness of reconciliation itself from the manner in which God accomplishes it. Estranged humanity is reconciled to God; the Son himself does this, through his death and through the Cross. The Creator of those estranged from him submits to death for their reconciliation.\footnote{Chrysostom, \emph{Homilies on Colossians} III, exposition of Col 1:19--20.}

Jesus' conduct before Pilate can be understood in relation to this reconciliation. His servants do not fight to prevent his delivery, and the Passion proceeds. The Gospel recounts the trial of the one whose death the Epistle identifies as the means of peace. His refusal of rescue allows that self-offering to be accomplished.

The Introit's heavenly worshippers acclaim the Lamb as slain. The Collect asks for the restoration of human families wounded by sin. In the Secret, the Father is asked that the offered Christ give those peoples unity and peace. The same Christ is acclaimed for his sacrifice and offered for the reconciliation of those who still need its fruits.

The Epistle describes the recipients as former enemies at Colossians 1:21, immediately after the appointed passage. Their reconciliation requires a change in their relation to God. The prayer for divided human families asks for such a change among them as well. A community receiving this peace must acknowledge the wrongdoing that divides it and seek the conversion of those responsible, including its own members.

\subsection{What is reconciled in heaven?}
The Epistle's phrase about things in heaven requires explanation. Chrysostom asks how there can be war among heavenly beings when heaven is the model of obedience invoked in the Lord's Prayer. His answer does not make the holy angels sinners. The breach lies between heaven and sinful humanity: the angels oppose the dishonor humanity has done to their Lord.\footnote{Chrysostom, \emph{Homilies on Colossians} III, discussion beginning with the question about ``things in the heavens.''}

Christ brings humanity into heavenly company and removes the estrangement. Chrysostom points to the angels' joyful appearance among human beings and to humanity's reception in heaven. Earth is reconciled both internally and with heaven; heaven receives the human beings formerly alienated from it. The angels' joy accompanies the peace the Son has made.

This explains why Chrysostom turns against reliance upon angels as substitutes for Christ. Their ministry is real, but they are themselves servants within the order the Son reconciles. To make them alternative saviors would mistake the beneficiaries and witnesses of his peace for its independent source. The Son's creative primacy and reconciling action thus answer one connected concern: no invisible power can offer a saving route that bypasses him.

The Preface recounts the priestly King's self-offering before the assembly gives thanks with the angels. Chrysostom's explanation of heavenly reconciliation helps explain this shared worship. Human beings estranged from heaven are received into its company through the Son's work. In this interpretation, the assembly's praise with the angels expresses that restored communion.

The universality of the claim must retain the Epistle's immediate context. Colossians 1:21--23 turns from the cosmic horizon to former enemies and their perseverance in faith. Chrysostom's explanation of heavenly reconciliation concerns holy angels and humanity, not universal pardon for every rebellious spiritual being. Nor does the scale of Christ's work imply that every human person already enjoys saving friendship with him. The kingdom's reach and the response of its subjects remain distinct.

\subsection{All kings, and the recognition still to come}
The Gradual declares a dominion embracing all nations and rulers. Bellarmine asks how such language can be true when willing acknowledgment is visibly incomplete. His commentary on Psalm 71:11 presents three ways the universal claim bears upon Christ. It can describe widespread conversion in the scriptural manner of speaking generally while allowing exceptions; it can look toward final judgment, when acknowledgment will be complete; and it can name the actual authority Christ already has over all, with or without their consent.\footnote{Bellarmine, \emph{Commentary on Psalms}, Ps 71, numbered comment 10 on biblical v. 11, O'Sullivan translation.}

The alternatives are not interchangeable. Conversion describes a change in the subject. Final acknowledgment describes an encounter with judgment that cannot be evaded. Present authority describes the King's right and power even where the subject resists. Confusing these claims produces two opposite errors: announcing that everyone already obeys willingly, or deciding that resistance proves Christ lacks universal authority. Bellarmine avoids both by explaining what each sense asserts.

Bellarmine also allows the prayer for the king's judgment to concern David and Solomon in its immediate historical reference. Solomon can receive a responsibility to govern justly; the request is intelligible within Israel's royal life. Bellarmine then interprets the promises of endurance and universal dominion as fulfilled in Christ, since they exceed Solomon's reign. Augustine begins instead from Christ as true Solomon and peacemaker. Their disagreement concerns the immediate scope of the opening, while both find its fulfillment in Christ.\footnote{Bellarmine, \emph{Commentary on Psalms}, Ps 71, opening and comments 4--7; Augustine, \emph{Expositions of the Psalms}, Ps 71.1--2.}

The verses following universal homage explain why people turn to this King. He rescues the helpless. Bellarmine develops that rescue as liberation from the devil's captivity through forgiveness and restoration of grace. Human poverty is not a reason for the King to despise his subjects; it is the condition in which they discover their need of him. His account of deliverance can be read beside the Epistle's transfer from darkness: the Son's rule frees those held captive and restores their life with God.\footnote{Bellarmine, \emph{Commentary on Psalms}, Ps 71, comments 11--13, on the deliverance passage following biblical v. 11.}

At the offering, Psalm 2:8 names the ends of the earth as the extent of the Son's inheritance. Augustine relates the Son's asking to priestly intercession; Bellarmine relates the inheritance to the Son's identity and rightful possession. His authority extends to peoples who still need conversion. The petitions for their healing in the Collect and Secret therefore remain necessary even as the Mass acclaims his universal rule.

\subsection{The Church's peace and the feud at its door}
Chrysostom does not allow reconciliation to remain a grand picture remote from ordinary conduct. His homily turns to the Church's repeated exchange of peace. The greeting is directed to all. A person who maintains peace with one neighbor and a feud with another cannot treat partial agreement as the health of the whole body. Chrysostom compares such selective peace to a body in which some members are at rest while others remain disordered. The body is not healthy merely because the disorder is localized.\footnote{Chrysostom, \emph{Homilies on Colossians} III, practical application following the exposition of heavenly reconciliation.}

Chrysostom applies the comparison to the person who receives the Church's greeting of peace while deliberately preserving a feud. The Son came to reconcile heaven and earth, and the recipient refuses reconciliation with a fellow member of the Church. For Chrysostom, this conduct contradicts the peace being exchanged. His application makes the healing of a particular relationship part of the response to Christ's reconciling work.

The Communion names peace as the King's blessing. The Postcommunion asks that those already nourished by the sacrament and glorying in service under Christ's standard may reign with him in heaven. In this interpretation, their present service includes reconciliation with others in response to the peace they receive. Augustine's exposition of Psalm 28 describes peace in Christ amid continuing opposition. Perseverance in that peace can require truthfulness, patient correction and resistance to wrongdoing, while excluding hatred of the person who resists.

The Alleluia sings of the dominion that outlasts the unstable powers of Daniel's vision. Jerome explains its reception by the Son made man and its enduring possession by the saints. The Postcommunion asks for a share in that heavenly reign. The peace received and practiced amid present conflict reaches its fulfillment in lasting life with Christ.

\begin{fourSenses}
\item[Literal.] Colossians joins creation, headship and the Cross; the royal Psalms join dominion to justice; Jesus' trial discloses a kingship whose witness proceeds through his delivery; the prayers seek reconciliation and peace.
\item[Allegorical.] The one Son is Creator, incarnate head, priestly victim and King. His self-giving restores humanity to communion with God and heavenly worship, and gathers his Church from the nations.
\item[Moral.] Those who receive peace must oppose their own feuds and care for the helpless. Royal justice requires conversion and public responsibility; voluntary allegiance is more than unavoidable subjection to authority.
\item[Anagogical.] Universal authority will receive complete acknowledgment at judgment. Saving communion reaches the eternal kingdom, while present resistance and imperfect recognition neither exhaust nor annul Christ's reign.
\end{fourSenses}
